%TCIDATA{Version=5.50.0.2953}
%TCIDATA{LaTeXparent=0,0,EOMN.tex}
                      
%TCIDATA{ChildDefaults=chapter:1,page:1}


\section{Mappings between pairs of Dual Spaces}

Consider a {\em compact} map 
\begin{equation}
A:V\longrightarrow W
\end{equation}%
and its adjoint  
\begin{equation}
A^{\dagger }:W^{d}\longrightarrow V^{d}
\end{equation}%
where $\left\langle V^{d},V\right\rangle $ and $\left\langle
W^{d},W\right\rangle $ are dual space pairs and 
\begin{equation}
\left\langle A^{\dagger }{\bf w}^{d}|{\bf v}\right\rangle =\left\langle {\bf %
w}^{d}|A{\bf v}\right\rangle \quad \forall {\bf v\in }V
\end{equation}%
We then have the following diagram 
\begin{equation}
\frame{$%
\begin{array}{ccc}
V & \stackrel{A}{%
%TCIMACRO{\underset{A^{-1}}{{\LARGE \rightleftarrows }}}%
%BeginExpansion
\mathrel{\mathop{{\LARGE \rightleftarrows }}\limits_{A^{-1}}}%
%EndExpansion
} & W \\ 
&  &  \\ 
V^{d} & \stackrel{A^{\dagger }}{%
%TCIMACRO{\underset{A^{-1\dagger }}{{\LARGE \rightleftarrows }}}%
%BeginExpansion
\mathrel{\mathop{{\LARGE \rightleftarrows }}\limits_{A^{-1\dagger }}}%
%EndExpansion
} & W^{d}%
\end{array}%
$}
\end{equation}%
As the maps are compact their domain and range are closed subspaces and we
can {\it wlog }take 
\begin{eqnarray}
V &=&%
%TCIMACRO{\func{domain}}%
%BeginExpansion
\mathop{\rm domain}%
%EndExpansion
\left\{ A\right\}   \nonumber \\
W^{d} &=&%
%TCIMACRO{\func{domain}}%
%BeginExpansion
\mathop{\rm domain}%
%EndExpansion
\left\{ A^{\dagger }\right\} 
\end{eqnarray}%
Further we can embed $V\hookrightarrow $ $V^{dd}$ by ${\bf v\rightarrow \hat{%
v}}$ 
\begin{equation}
\left\langle {\bf \hat{v}}|{\bf v}^{d}\right\rangle =\left\langle {\bf v}%
^{d}|{\bf v}\right\rangle \quad \forall {\bf v}^{d}\in V^{d}\text{ }
\end{equation}%
so that $V^{dd}\supseteq V$ with equality always occurring for finite
dimensional spaces. The same holds true for the $W$ spaces.

Let

\begin{itemize}
\item $\left\{ \left| k_{j}\right\rangle ;1\leq j\leq m_{1}\right\} $ be a
basis for $\ker A$ and $\left\{ \left\langle k_{j}^{d}\right| ;1\leq j\leq
m_{1}\right\} $ be a biothorgonal dual basis of $\left\{ \ker A\right\} ^{d}$

\item $\left\{ \left| s_{j}\right\rangle ;1\leq j\leq m_{2}\right\} $ be a
basis for $\left\{ \ker A\right\} ^{c}$ and $\left\{ \left\langle
s_{j}^{d}\right| ;1\leq j\leq m_{2}\right\} $ be a biothorgonal dual basis
of $\left\{ \ker A\right\} ^{cd}$

\item $\left\{ \left| \tilde{s}_{j}\right\rangle ;1\leq j\leq m_{2}\right\} $
be a basis for the $%
%TCIMACRO{\func{range}}%
%BeginExpansion
\mathop{\rm range}%
%EndExpansion
A$ and $\left\{ \left\langle \tilde{s}_{j}^{d}\right| ;1\leq j\leq
m_{2}\right\} $ be a biothorgonal dual basis of $\left\{ 
%TCIMACRO{\func{range}}%
%BeginExpansion
\mathop{\rm range}%
%EndExpansion
A\right\} ^{d}$

\item $\left\{ \left| l_{j}\right\rangle ;1\leq j\leq m_{3}\right\} $ be a
basis for $\left\{ 
%TCIMACRO{\func{range}}%
%BeginExpansion
\mathop{\rm range}%
%EndExpansion
A\right\} ^{c}$ and $\left\{ \left\langle l_{j}^{d}\right| ;1\leq j\leq
m_{3}\right\} $ be a biothorgonal dual basis of $\left\{ 
%TCIMACRO{\func{range}}%
%BeginExpansion
\mathop{\rm range}%
%EndExpansion
A\right\} ^{cd}$.
\end{itemize}

Note that 
\begin{eqnarray}
V &=&\ker A\oplus \left\{ \ker A\right\} ^{c}  \nonumber \\
V^{d} &=&\left\{ \ker A\right\} ^{d}\oplus \left\{ \ker A\right\} ^{cd}
\end{eqnarray}

and 
\begin{eqnarray}
W &=&%
%TCIMACRO{\func{range}}%
%BeginExpansion
\mathop{\rm range}%
%EndExpansion
A\oplus \left\{ 
%TCIMACRO{\func{range}}%
%BeginExpansion
\mathop{\rm range}%
%EndExpansion
A\right\} ^{c}  \nonumber \\
W^{d} &=&\left\{ 
%TCIMACRO{\func{range}}%
%BeginExpansion
\mathop{\rm range}%
%EndExpansion
A\right\} ^{d}\oplus \left\{ 
%TCIMACRO{\func{range}}%
%BeginExpansion
\mathop{\rm range}%
%EndExpansion
A\right\} ^{cd}
\end{eqnarray}%
Thus 
\begin{eqnarray}
\dim V &=&m_{1}+m_{2}  \nonumber \\
\dim W &=&m_{2}+m_{3}  \nonumber \\
\dim V-\dim W &=&m_{1}-m_{3}
\end{eqnarray}%
One can observe that $\dim \left\{ 
%TCIMACRO{\func{range}}%
%BeginExpansion
\mathop{\rm range}%
%EndExpansion
A\right\} =\dim \left\{ \ker A\right\} ^{c}=m_{2}$. Associated with these
spaces are the projectors 
\begin{eqnarray}
P_{\ker A} &=&\sum_{1\leq j\leq m_{1}}\left| k_{j}\right\rangle \left\langle
k_{j}^{d}\right|  \nonumber \\
P_{\left\{ \ker A\right\} ^{c}} &=&\sum_{1\leq j\leq m_{2}}\left|
s_{j}\right\rangle \left\langle s_{j}^{d}\right|
\end{eqnarray}%
which have the properties%
\begin{eqnarray}
P_{\ker A}^{2} &=&P_{\ker A}  \nonumber \\
P_{\left\{ \ker A\right\} ^{c}}^{2} &=&P_{\left\{ \ker A\right\} ^{c}} 
\nonumber \\
P_{\left\{ \ker A\right\} ^{c}}P_{\ker A} &=&P_{\ker A}P_{\left\{ \ker
A\right\} ^{c}}=0  \nonumber \\
I_{V} &=&P_{\ker A}+P_{\left\{ \ker A\right\} ^{c}}
\end{eqnarray}%
and 
\begin{eqnarray}
P_{%
%TCIMACRO{\func{range}}%
%BeginExpansion
\mathop{\rm range}%
%EndExpansion
A} &=&\sum_{1\leq j\leq m_{2}}\left| \tilde{s}_{j}\right\rangle \left\langle 
\tilde{s}_{j}^{d}\right|  \nonumber \\
P_{\left\{ 
%TCIMACRO{\func{range}}%
%BeginExpansion
\mathop{\rm range}%
%EndExpansion
A\right\} ^{c}} &=&\sum_{1\leq j\leq m_{3}}\left| l_{j}\right\rangle
\left\langle l_{j}^{d}\right|
\end{eqnarray}%
which have the properties%
\begin{eqnarray}
P_{%
%TCIMACRO{\func{range}}%
%BeginExpansion
\mathop{\rm range}%
%EndExpansion
A}^{2} &=&P_{%
%TCIMACRO{\func{range}}%
%BeginExpansion
\mathop{\rm range}%
%EndExpansion
A}  \nonumber \\
P_{\left\{ 
%TCIMACRO{\func{range}}%
%BeginExpansion
\mathop{\rm range}%
%EndExpansion
A\right\} ^{c}}^{2} &=&P_{\left\{ 
%TCIMACRO{\func{range}}%
%BeginExpansion
\mathop{\rm range}%
%EndExpansion
A\right\} ^{c}}  \nonumber \\
P_{\left\{ 
%TCIMACRO{\func{range}}%
%BeginExpansion
\mathop{\rm range}%
%EndExpansion
A\right\} ^{c}}P_{%
%TCIMACRO{\func{range}}%
%BeginExpansion
\mathop{\rm range}%
%EndExpansion
A} &=&P_{%
%TCIMACRO{\func{range}}%
%BeginExpansion
\mathop{\rm range}%
%EndExpansion
A}P_{\left\{ 
%TCIMACRO{\func{range}}%
%BeginExpansion
\mathop{\rm range}%
%EndExpansion
A\right\} ^{c}}=0  \nonumber \\
I_{W} &=&P_{%
%TCIMACRO{\func{range}}%
%BeginExpansion
\mathop{\rm range}%
%EndExpansion
A}+P_{\left\{ 
%TCIMACRO{\func{range}}%
%BeginExpansion
\mathop{\rm range}%
%EndExpansion
A\right\} ^{c}}
\end{eqnarray}

If $A$ is compact and these bases are canonical for $A$ we obtain the
following expansions for $A$ 
\begin{equation}
A=\sum_{1\leq j\leq m_{2}}\alpha _{j}\left\vert \tilde{s}_{j}\right\rangle
\left\langle s_{j}^{d}\right\vert ;\quad \alpha _{j}>0
\end{equation}%
It is easy to see that $A$ has the properties 
\begin{eqnarray}
AP_{\left\{ \ker A\right\} ^{c}} &=&A  \nonumber \\
P_{%
%TCIMACRO{\func{range}}%
%BeginExpansion
\mathop{\rm range}%
%EndExpansion
A}A &=&A  \nonumber \\
P_{%
%TCIMACRO{\func{range}}%
%BeginExpansion
\mathop{\rm range}%
%EndExpansion
A}AP_{\left\{ \ker A\right\} ^{c}} &=&A
\end{eqnarray}%
By utilizing the canonical expansions of $A$ we can define a unique two
sided generalized inverse 
\begin{equation}
A^{-1}=\sum_{1\leq j\leq m_{2}}\alpha _{j}^{-1}\left\vert s_{j}\right\rangle
\left\langle \tilde{s}_{j}^{d}\right\vert  \label{Ainv}
\end{equation}%
which has the property 
\begin{equation}
AA^{-1}=\sum_{1\leq j\leq m_{2}}\left\vert \tilde{s}_{j}\right\rangle
\left\langle \tilde{s}_{j}^{d}\right\vert =P_{%
%TCIMACRO{\func{range}}%
%BeginExpansion
\mathop{\rm range}%
%EndExpansion
A}
\end{equation}%
and 
\begin{equation}
A^{-1}A=\sum_{1\leq j\leq m_{2}}\left\vert s_{j}\right\rangle \left\langle
s_{j}^{d}\right\vert =P_{\left\{ \ker A\right\} ^{c}}
\end{equation}%
and\ 
\begin{eqnarray}
A^{-1}P_{%
%TCIMACRO{\func{range}}%
%BeginExpansion
\mathop{\rm range}%
%EndExpansion
A} &=&A^{-1}  \nonumber \\
P_{\left\{ \ker A\right\} ^{c}}A^{-1} &=&A^{-1}  \nonumber \\
P_{\left\{ \ker A\right\} ^{c}}A^{-1}P_{%
%TCIMACRO{\func{range}}%
%BeginExpansion
\mathop{\rm range}%
%EndExpansion
A} &=&A^{-1}
\end{eqnarray}%
The definition in eq. $\left( \ref{Ainv}\right) $ can be used as a {\em core}
for left and right inverses of $A$ as is shown in the following. Let 
\begin{eqnarray}
R &:&\left\{ 
%TCIMACRO{\func{range}}%
%BeginExpansion
\mathop{\rm range}%
%EndExpansion
A\right\} ^{c}\longrightarrow \ker A  \nonumber \\
R &=&\sum R_{ij}\left\vert k_{i}\right\rangle \left\langle
l_{j}^{d}\right\vert
\end{eqnarray}%
and define 
\begin{equation}
A_{R}^{-1}=A^{-1}+R
\end{equation}%
then 
\begin{equation}
AA_{R}^{-1}=\sum_{1\leq j\leq m_{2}}\alpha _{j}\left\vert \tilde{s}%
_{j}\right\rangle \left\langle s_{j}^{d}\right\vert \left( \sum_{1\leq j\leq
m_{2}}\alpha _{j}^{-1}\left\vert s_{j}\right\rangle \left\langle \tilde{s}%
_{j}^{d}\right\vert +\sum_{\Sb 1\leq i\leq m_{1}  \\ 1\leq j\leq m_{3} 
\endSb }R_{ij}\left\vert k_{i}\right\rangle \left\langle
l_{j}^{d}\right\vert \right) =P_{%
%TCIMACRO{\func{range}}%
%BeginExpansion
\mathop{\rm range}%
%EndExpansion
A}
\end{equation}%
and 
\begin{eqnarray}
L &:&\ker A\longrightarrow \left\{ 
%TCIMACRO{\func{range}}%
%BeginExpansion
\mathop{\rm range}%
%EndExpansion
A\right\} ^{c}  \nonumber \\
L &=&\sum\Sb 1\leq i\leq m_{3}  \\ 1\leq j\leq m_{1}  \endSb %
L_{ij}\left\vert k_{i}\right\rangle \left\langle l_{j}^{d}\right\vert
\end{eqnarray}%
and define 
\begin{equation}
A_{L}^{-1}=A^{-1}+L
\end{equation}%
then 
\begin{equation}
A_{L}^{-1}A=\left( \sum_{1\leq j\leq m_{2}}\alpha _{j}^{-1}\left\vert
s_{j}\right\rangle \left\langle \tilde{s}_{j}^{d}\right\vert +\sum\Sb 1\leq
i\leq m_{3}  \\ 1\leq j\leq m_{1}  \endSb L_{ij}\left\vert
k_{i}\right\rangle \left\langle l_{j}^{d}\right\vert \right) \sum_{1\leq
j\leq m_{2}}\alpha _{j}\left\vert \tilde{s}_{j}\right\rangle \left\langle
s_{j}^{d}\right\vert =P_{\left\{ \ker A\right\} ^{c}}
\end{equation}%
The preceding relationships are an abstract way of defining left and right
inverses of $A$, based on a {\em unique} core generalized inverse.

The adjoint $A^{\dagger }$ of $A$ can be expanded as 
\begin{equation}
A^{\dagger }=\sum_{1\leq j\leq m_{2}}\alpha _{j}\left\vert \tilde{s}%
_{j}^{d}\right\rangle \left\langle \tilde{s}_{j}\right\vert ;\quad \alpha
_{j}>0
\end{equation}%
We define the orthogonal compliment of a subspace $U\subseteq V$ as a
subspace $U^{\bot }$ contained in the dual space $V^{d}$ as 
\begin{equation}
U^{\bot }=\left\{ {\bf u}^{d};\backepsilon \left\langle {\bf u}^{d}|{\bf u}%
\right\rangle =0\quad \forall {\bf u\in }U\right\}
\end{equation}%
and note that 
\begin{equation}
U^{\bot \bot }\supseteq U
\end{equation}%
Further one has 
\begin{equation}
U^{cd}=U^{\bot }
\end{equation}%
Finally we note a relationship between $\ker A$ and $%
%TCIMACRO{\func{range}}%
%BeginExpansion
\mathop{\rm range}%
%EndExpansion
A^{\dagger }$. If ${\bf v}\in \ker A$ then%
\begin{equation}
\left\langle {\bf w}^{d}|A{\bf v}\right\rangle =0\,\forall {\bf w}%
^{d}\Rightarrow \left\langle A^{\dagger }{\bf w}^{d}|{\bf v}\right\rangle
=0\,\forall {\bf w}^{d}\Rightarrow {\bf v}\in \left\{ 
%TCIMACRO{\func{range}}%
%BeginExpansion
\mathop{\rm range}%
%EndExpansion
A^{\dagger }\right\} ^{\bot }\supseteq \ker A
\end{equation}%
Similarly 
\begin{equation}
\left\langle A^{\dagger }{\bf w}^{d}|{\bf v}\right\rangle =0\,\forall {\bf v}%
\Rightarrow \left\langle {\bf w}^{d}|A{\bf v}\right\rangle =0\,\forall {\bf v%
}\Rightarrow {\bf w}^{d}\in \left\{ 
%TCIMACRO{\func{range}}%
%BeginExpansion
\mathop{\rm range}%
%EndExpansion
A\right\} ^{\bot }\supseteq \ker A^{\dagger }
\end{equation}%
$\lceil $Recall $\left\langle A^{\dagger }{\bf w}^{d}|{\bf v}\right\rangle
=\left\langle {\bf w}^{d}|A{\bf v}\right\rangle \quad \forall {\bf v\rfloor }%
\ \ $

In many applications we have the situation that $%
%TCIMACRO{\func{domain}}%
%BeginExpansion
\mathop{\rm domain}%
%EndExpansion
\left\{ A\right\} \supseteq 
%TCIMACRO{\func{range}}%
%BeginExpansion
\mathop{\rm range}%
%EndExpansion
\left\{ A\right\} $ (equality occurs if $A$ is $1-1$) so that 
\begin{eqnarray}
%TCIMACRO{\func{domain}}%
%BeginExpansion
\mathop{\rm domain}%
%EndExpansion
\left\{ A\right\} &=&\ker \left\{ A\right\} \oplus \ker \left\{ A\right\}
^{c}  \nonumber \\
%TCIMACRO{\func{range}}%
%BeginExpansion
\mathop{\rm range}%
%EndExpansion
\left\{ A\right\} &\cong &\ker \left\{ A\right\} ^{c}
\end{eqnarray}%
thus 
\begin{equation}
%TCIMACRO{\func{range}}%
%BeginExpansion
\mathop{\rm range}%
%EndExpansion
\left\{ A\right\} ^{c}\cong \ker \left\{ A\right\}
\end{equation}%
and 
\begin{equation}
%TCIMACRO{\func{domain}}%
%BeginExpansion
\mathop{\rm domain}%
%EndExpansion
\left\{ A\right\} =%
%TCIMACRO{\func{range}}%
%BeginExpansion
\mathop{\rm range}%
%EndExpansion
\left\{ A\right\} ^{c}\oplus 
%TCIMACRO{\func{range}}%
%BeginExpansion
\mathop{\rm range}%
%EndExpansion
\left\{ A\right\}
\end{equation}%
but note that $\ker \left\{ A\right\} \neq 
%TCIMACRO{\func{range}}%
%BeginExpansion
\mathop{\rm range}%
%EndExpansion
\left\{ A\right\} ^{c}$ unless $\ker \left\{ A\right\} ^{c}$ is an invariant
subspace of $A$, which it is if $A$ is self adjoint.

Another common simplification occurs when $V$ \ and $W$ are self dual i.e. $%
V=V^{d}$ and $W=W^{d}$ and are thus Hilbert spaces. When both of the
preceding simplifications occur the projectors become orthogonal projectors
and all the vectors belong to the same space. Though note that the dual
basis and the normal basis are still in general distinct, unless $A$ is self
adjoint.

\subsection{Trace}

\begin{equation}
A=\sum_{1\leq j\leq m_{2}}\alpha _{j}\left\vert \tilde{s}_{j}\right\rangle
\left\langle s_{j}^{d}\right\vert ;\quad \alpha _{j}>0
\end{equation}%
\begin{equation}
A^{\dagger }=\sum_{1\leq j\leq m_{2}}\alpha _{j}\left\vert \tilde{s}%
_{j}^{d}\right\rangle \left\langle \tilde{s}_{j}\right\vert ;\quad \alpha
_{j}>0
\end{equation}%
\begin{eqnarray*}
AA^{\dagger } &=&\sum_{1\leq j\leq m_{2}}\alpha _{j}\left\vert \tilde{s}%
_{j}\right\rangle \left\langle s_{j}^{d}\right\vert \sum_{1\leq j\leq
m_{2}}\alpha _{j}\left\vert \tilde{s}_{j}^{d}\right\rangle \left\langle 
\tilde{s}_{j}\right\vert = \\
A^{\dagger }A &=&\sum_{1\leq j\leq m_{2}}\alpha _{j}\left\vert \tilde{s}%
_{j}^{d}\right\rangle \left\langle \tilde{s}_{j}\right\vert \sum_{1\leq
j\leq m_{2}}\alpha _{j}\left\vert \tilde{s}_{j}\right\rangle \left\langle
s_{j}^{d}\right\vert =
\end{eqnarray*}

\subsection{Linear Equations}

We can now reexamine the linear equation 
\begin{equation}
A\left| {\bf v}\right\rangle =\left| {\bf w}\right\rangle {\bf ;\quad \left|
v\right\rangle \in V}^{s},\,\left| {\bf w}\right\rangle {\bf \in W}%
^{r};\quad r=m_{2}+m_{3};\quad s=m_{2}+m_{1}  \label{lineq}
\end{equation}%
This equation can be written as%
\begin{eqnarray}
&&\sum_{1\leq j\leq m_{2}}\alpha _{j}\left| \tilde{s}_{j}\right\rangle
\left\langle s_{j}^{d}\right| {\bf v}\rangle =\left| {\bf w}\right\rangle
\,\,{\bf \Rightarrow \,}P_{\left\{ \ker A\right\} ^{c}}{\bf \,}\left| {\bf v}%
\right\rangle =A_{L}^{-1}\left| {\bf w}\right\rangle  \nonumber \\
&\equiv &\left( \sum_{1\leq j\leq m_{2}}\alpha _{j}^{-1}\left|
s_{j}\right\rangle \left\langle \tilde{s}_{j}^{d}\right| +\sum\Sb 1\leq
i\leq m_{3}  \\ 1\leq j\leq m_{1}  \endSb L_{ij}\left| k_{i}\right\rangle
\left\langle l_{j}^{d}\right| \right) \sum_{1\leq j\leq m_{2}}\alpha
_{j}\left| \tilde{s}_{j}\right\rangle \left\langle s_{j}^{d}\right| {\bf v}%
\rangle =P_{\left\{ \ker A\right\} ^{c}}\left| {\bf v}\right\rangle 
\nonumber \\
&{\bf =}&\left( \sum_{1\leq j\leq m_{2}}\alpha _{j}^{-1}\left|
s_{j}\right\rangle \left\langle \tilde{s}_{j}^{d}\right| +\sum\Sb 1\leq
i\leq m_{3}  \\ 1\leq j\leq m_{1}  \endSb L_{ij}\left| k_{i}\right\rangle
\left\langle l_{j}^{d}\right| \right) \left| {\bf w}\right\rangle
\end{eqnarray}%
i.e. 
\begin{equation}
P_{\left\{ \ker A\right\} ^{c}}\left| {\bf v}\right\rangle =\sum_{1\leq
j\leq m_{2}}\alpha _{j}^{-1}\left| s_{j}\right\rangle \left\langle \tilde{s}%
_{j}^{d}\right| {\bf w\rangle +}\sum\Sb 1\leq i\leq m_{3}  \\ 1\leq j\leq
m_{1}  \endSb L_{ij}\left| k_{i}\right\rangle \left\langle l_{j}^{d}\right| 
{\bf w\rangle }
\end{equation}%
which can only be true if the consistency constraint 
\begin{equation}
\left\langle l_{j}^{d}\right| {\bf w}\rangle =0;\quad 1\leq i\leq
m_{3}\Longleftrightarrow P_{\left\{ 
%TCIMACRO{\func{range}}%
%BeginExpansion
\mathop{\rm range}%
%EndExpansion
A\right\} ^{c}}\left| {\bf w}\right\rangle =0
\end{equation}%
holds. The remaining part of the vector $\left| {\bf v}\right\rangle $, i.e. 
$P_{\left\{ \ker A\right\} }\left| {\bf v}\right\rangle $ is arbitrary
unless some further conditions are invoked. Thus the solution to eq. $\left( %
\ref{lineq}\right) $ is 
\begin{equation}
\left| {\bf v}\right\rangle =A^{-1}\left| {\bf w}\right\rangle +\left| {\bf k%
}\right\rangle ;\,\,\left| {\bf k}\right\rangle \in \ker A
\end{equation}%
as long as the condition 
\begin{equation}
P_{\left\{ 
%TCIMACRO{\func{range}}%
%BeginExpansion
\mathop{\rm range}%
%EndExpansion
A\right\} ^{c}}\left| {\bf w}\right\rangle =0
\end{equation}%
is satisfied.

\subsection{Matrix Representations}

Using the preceding biorthogonal basis we obtain 
\begin{eqnarray}
&&A\left| \left( {\bf s,k}\right) \right\rangle =\left| \left( {\bf \tilde{s}%
,l}\right) \right\rangle \left( 
\begin{array}{cc}
\begin{array}{ccc}
\alpha _{1} &  &  \\ 
& \ddots &  \\ 
&  & \alpha _{s}%
\end{array}
& 0 \\ 
0 & 0%
\end{array}%
\right) ;\quad \text{ }\left( m_{1}+m_{3}\right) \times \left(
m_{1}+m_{2}\right) \text{ matrix}  \nonumber \\
&{\Bbb \Rightarrow }&\left\langle \left( {\bf \tilde{s}}^{d},{\bf l}%
^{d}\right) |A\left( {\bf s,k}\right) \right\rangle =\left\langle \left( 
{\bf \tilde{s}}^{d},{\bf l}^{d}\right) |\left( {\bf s,k}\right)
\right\rangle \left( 
\begin{array}{cc}
\begin{array}{ccc}
\alpha _{1} &  &  \\ 
& \ddots &  \\ 
&  & \alpha _{s}%
\end{array}
& 0 \\ 
0 & 0%
\end{array}%
\right) =  \nonumber \\
&&\qquad \qquad \qquad \qquad \quad {\Bbb I}_{\left( m_{1}+m_{3}\right)
}\left( 
\begin{array}{cc}
\begin{array}{ccc}
\alpha _{1} &  &  \\ 
& \ddots &  \\ 
&  & \alpha _{s}%
\end{array}
& 0 \\ 
0 & 0%
\end{array}%
\right) =\left( 
\begin{array}{cc}
\begin{array}{ccc}
\alpha _{1} &  &  \\ 
& \ddots &  \\ 
&  & \alpha _{s}%
\end{array}
& 0 \\ 
0 & 0%
\end{array}%
\right)
\end{eqnarray}%
and 
\begin{equation}
A^{-1}\left( \left| {\bf \tilde{s}}\right\rangle {\bf ,}\left| {\bf l}%
\right\rangle \right) =\left( \left| {\bf s}\right\rangle {\bf ,}\left| {\bf %
k}\right\rangle \right) \left( 
\begin{array}{cc}
\begin{array}{ccc}
\alpha _{1}^{-1} &  &  \\ 
& \ddots &  \\ 
&  & \alpha _{s}^{-1}%
\end{array}
& 0 \\ 
0 & 0%
\end{array}%
\right) ;\quad \text{ }\left( m_{1}+m_{2}\right) \times \left(
m_{1}+m_{3}\right) \text{ matrix}
\end{equation}%
Hence 
\begin{equation}
\left( 
\begin{array}{cc}
\begin{array}{ccc}
\alpha _{1}^{-1} &  &  \\ 
& \ddots &  \\ 
&  & \alpha _{s}^{-1}%
\end{array}
& 0 \\ 
0 & 0%
\end{array}%
\right) \left( 
\begin{array}{cc}
\begin{array}{ccc}
\alpha _{1} &  &  \\ 
& \ddots &  \\ 
&  & \alpha _{s}%
\end{array}
& 0 \\ 
0 & 0%
\end{array}%
\right) =\left( 
\begin{array}{cc}
\begin{array}{ccc}
1 &  &  \\ 
& \ddots &  \\ 
&  & 1%
\end{array}
& 0 \\ 
0 & 0%
\end{array}%
\right) ={\Bbb P}_{\left\{ \ker A\right\} ^{c}}
\end{equation}%
which is $\left( m_{1}+m_{2}\right) \times \left( m_{1}+m_{2}\right) $
matrix, and 
\begin{equation}
\left( 
\begin{array}{cc}
\begin{array}{ccc}
\alpha _{1} &  &  \\ 
& \ddots &  \\ 
&  & \alpha _{s}%
\end{array}
& 0 \\ 
0 & 0%
\end{array}%
\right) \left( 
\begin{array}{cc}
\begin{array}{ccc}
\alpha _{1}^{-1} &  &  \\ 
& \ddots &  \\ 
&  & \alpha _{s}^{-1}%
\end{array}
& 0 \\ 
0 & 0%
\end{array}%
\right) =\left( 
\begin{array}{cc}
\begin{array}{ccc}
1 &  &  \\ 
& \ddots &  \\ 
&  & 1%
\end{array}
& 0 \\ 
0 & 0%
\end{array}%
\right) ={\Bbb P}_{%
%TCIMACRO{\func{range}}%
%BeginExpansion
\mathop{\rm range}%
%EndExpansion
A}
\end{equation}%
which is $\left( m_{1}+m_{3}\right) \times \left( m_{1}+m_{3}\right) $
matrix.

If we choose other bases related to the canonical basis of $A$ by 
\begin{eqnarray}
&&\qquad \quad \left| \left( {\bf s}^{\prime }{\bf ,k}^{\prime }\right)
\right\rangle =S\left| \left( {\bf s,k}\right) \right\rangle =\left| \left( 
{\bf s,k}\right) \right\rangle {\Bbb S}  \nonumber \\
{\Bbb S\,} &{\Bbb =}&\left( 
\begin{array}{cc}
{\Bbb S}_{11} & {\Bbb S}_{12} \\ 
{\Bbb S}_{21} & {\Bbb S}_{22}%
\end{array}%
\right) =\left( 
\begin{array}{cc}
\left\langle {\bf s}^{d}|S{\bf s}\right\rangle & \left\langle {\bf s}^{d}|S%
{\bf k}\right\rangle \\ 
\left\langle {\bf k}^{d}|S{\bf s}\right\rangle & \left\langle {\bf k}^{d}|S%
{\bf k}\right\rangle%
\end{array}%
\right)  \nonumber \\
&&\qquad \quad \,\left| \left( {\bf \tilde{s}}^{\prime }{\bf ,l}^{\prime
}\right) \right\rangle =T\left| \left( {\bf \tilde{s},l}\right)
\right\rangle =\left| \left( {\bf \tilde{s},l}\right) \right\rangle {\Bbb T}
\nonumber \\
{\Bbb T}{\Bbb \,} &{\Bbb =}&\left( 
\begin{array}{cc}
{\Bbb T}_{11} & {\Bbb T}_{12} \\ 
{\Bbb T}_{21} & {\Bbb T}_{22}%
\end{array}%
\right) =\left( 
\begin{array}{cc}
\left\langle {\bf \tilde{s}}^{d}|T{\bf \tilde{s}}\right\rangle & 
\left\langle {\bf \tilde{s}}^{d}|T{\bf l}\right\rangle \\ 
\left\langle {\bf l}^{d}|T{\bf \tilde{s}}\right\rangle & \left\langle {\bf l}%
^{d}|T{\bf l}\right\rangle%
\end{array}%
\right)  \nonumber \\
&&\qquad \quad \Rightarrow \left| \left( {\bf \tilde{s},l}\right)
\right\rangle =\left| \left( {\bf \tilde{s}}^{\prime }{\bf ,l}^{\prime
}\right) \right\rangle {\Bbb T}^{-1}
\end{eqnarray}%
we obtain%
\begin{equation}
A\left| \left( {\bf s}^{\prime }{\bf ,k}^{\prime }\right) \right\rangle
=\left| \left( {\bf \tilde{s}}^{\prime }{\bf ,l}^{\prime }\right)
\right\rangle {\Bbb A}
\end{equation}%
where ${\Bbb A}$ is given by 
\begin{equation}
A\left| \left( {\bf s}^{\prime }{\bf ,k}^{\prime }\right) \right\rangle
=A\left| \left( {\bf s,k}\right) \right\rangle {\Bbb S=\,}\left| \left( {\bf 
\tilde{s},l}\right) \right\rangle {\Bbb \alpha S=\,}\left| \left( {\bf 
\tilde{s}}^{\prime }{\bf ,l}^{\prime }\right) \right\rangle {\Bbb T}^{-1}%
{\Bbb \alpha S}
\end{equation}
\begin{eqnarray}
A\left| \left( {\bf s}^{\prime }{\bf ,k}^{\prime }\right) \right\rangle
&=&\left| \left( {\bf \tilde{s}}^{\prime }{\bf ,l}^{\prime }\right)
\right\rangle \left( 
\begin{array}{cc}
\left( {\Bbb T}^{-1}\right) _{11} & \left( {\Bbb T}^{-1}\right) _{12} \\ 
\left( {\Bbb T}^{-1}\right) _{21} & \left( {\Bbb T}^{-1}\right) _{22}%
\end{array}%
\right) \left( 
\begin{array}{cc}
\begin{array}{ccc}
\alpha _{1} &  &  \\ 
& \ddots &  \\ 
&  & \alpha _{s}%
\end{array}
& 0 \\ 
0 & 0%
\end{array}%
\right) \left( 
\begin{array}{cc}
{\Bbb S}_{11} & {\Bbb S}_{12} \\ 
{\Bbb S}_{21} & {\Bbb S}_{22}%
\end{array}%
\right)  \nonumber \\
&=&\left| \left( {\bf \tilde{s}}^{\prime }{\bf ,l}^{\prime }\right)
\right\rangle \left( 
\begin{array}{cc}
\left( {\Bbb T}^{-1}\right) _{11} & \left( {\Bbb T}^{-1}\right) _{12} \\ 
\left( {\Bbb T}^{-1}\right) _{21} & \left( {\Bbb T}^{-1}\right) _{22}%
\end{array}%
\right) \left( 
\begin{array}{cc}
\alpha {\Bbb S}_{11} & \alpha {\Bbb S}_{12} \\ 
0 & 0%
\end{array}%
\right)  \nonumber \\
&=&\left| \left( {\bf \tilde{s}}^{\prime }{\bf ,l}^{\prime }\right)
\right\rangle \left( 
\begin{array}{cc}
\left( {\Bbb T}^{-1}\right) _{11}\alpha {\Bbb S}_{11} & \left( {\Bbb T}%
^{-1}\right) _{11}\alpha {\Bbb S}_{12} \\ 
\left( {\Bbb T}^{-1}\right) _{21}\alpha {\Bbb S}_{11} & \left( {\Bbb T}%
^{-1}\right) _{21}\alpha {\Bbb S}_{12}%
\end{array}%
\right) =\left| \left( {\bf \tilde{s}}^{\prime }{\bf ,l}^{\prime }\right)
\right\rangle {\Bbb A}
\end{eqnarray}%
matrix

Let ${\Bbb A}$ be a $r\times s$ matrix then its left and right inverses have
the property

\begin{eqnarray}
{\Bbb A}_{L}^{-1}{\Bbb A} &=&{\Bbb P}_{\left\{ \ker A\right\} ^{c}}\leq 
{\Bbb I}_{s}  \nonumber \\
{\Bbb AA}_{R}^{-1} &=&{\Bbb P}_{%
%TCIMACRO{\func{range}}%
%BeginExpansion
\mathop{\rm range}%
%EndExpansion
A}\leq {\Bbb I}_{r}
\end{eqnarray}%
and can be defined as matrix representations of the abstract operators
defined previously using the core generalized inverse.

\section{Solution of Nonlinear Singular Equations.}

Consider the following equation for $x$ 
\begin{equation}
{\cal A}\left( x\right) \left| b\left( x\right) \right\rangle =\left|
c\left( x\right) \right\rangle ;\quad \left| b\left( x\right) \right\rangle
\in V;\,\left| c\left( x\right) \right\rangle \in W  \label{nl}
\end{equation}%
where ${\cal A}\left( x\right) $ is a compact linear operator and $\left|
b\left( x\right) \right\rangle $ and $\left| c\left( x\right) \right\rangle $
are vectors that depends on $x.$ Let ${\cal A}\left( x\right) $ be singular
with a kernel, $\ker {\cal A}\left( x\right) $, and have the expansion 
\begin{equation}
{\cal A}\left( x\right) =\sum_{1\leq j\leq m_{2}}\alpha _{j}\left| \tilde{s}%
_{j}\right\rangle \left\langle s_{j}^{d}\right| ;\quad \alpha _{j}>0
\end{equation}%
The equation $\left( \ref{nl}\right) $ only has solution if $\left| c\left(
x\right) \right\rangle $ is in the range of ${\cal A}\left( x\right) $,
which gives the condition 
\begin{equation}
P_{\left\{ 
%TCIMACRO{\func{range}}%
%BeginExpansion
\mathop{\rm range}%
%EndExpansion
{\cal A}\left( x\right) \right\} ^{c}}\left| c\left( x\right) \right\rangle
=0
\end{equation}%
If this condition is enforced one can form the special solution 
\begin{equation}
{\frak s}\left( x\right) =\sum_{1\leq j\leq m_{2}}\alpha _{j}^{-1}\left|
s_{j}\right\rangle \left\langle \tilde{s}_{j}^{d}\right| \left. c\left(
x\right) \right\rangle
\end{equation}%
by using the generalized inverse ${\cal A}\left( x\right) ^{-1}$ of ${\cal A}%
\left( x\right) $ defined as 
\begin{equation}
{\cal A}\left( x\right) ^{-1}=\sum_{1\leq j\leq m_{2}}\alpha _{j}^{-1}\left|
s_{j}\right\rangle \left\langle \tilde{s}_{j}^{d}\right|
\end{equation}%
If ${\frak s}\left( x\right) $ is such a solution then the general solution
is given by 
\begin{equation}
{\frak g}\left( x\right) ={\frak s}\left( x\right) +\sum_{1\leq i\leq
m_{1}}\lambda _{i}\left| k_{i}\right\rangle
\end{equation}%
as 
\begin{equation}
\sum_{1\leq j\leq m_{2}}\alpha _{j^{\prime }}\left| \tilde{s}_{j^{\prime
}}\right\rangle \left\langle s_{j^{\prime }}^{d}\right| \left\{ \sum_{1\leq
j\leq m_{2}}\alpha _{j}^{-1}\left| s_{j}\right\rangle \left\langle \tilde{s}%
_{j}^{d}\right| \left. c\left( x\right) \right\rangle +\sum_{1\leq i\leq
m_{1}}\lambda _{i}\left| k_{i}\right\rangle \right\} =\left| c\left(
x\right) \right\rangle
\end{equation}%
Thus to summarize the solutions of eq. $\left( \ref{nl}\right) $ are $x\ni $ 
\begin{eqnarray}
&&b\left( x\right) ={\frak s}\left( x\right) +\sum_{1\leq i\leq
m_{1}}\lambda _{i}\left( x\right) \left| k_{i}\left( x\right) \right\rangle 
\nonumber \\
&&{\frak s}\left( x\right) =\sum_{1\leq j\leq m_{2}}\alpha _{j}^{-1}\left(
x\right) \left| s_{j}\left( x\right) \right\rangle \left\langle \tilde{s}%
_{j}^{d}\left( x\right) \right| \left. c\left( x\right) \right\rangle ={\cal %
A}\left( x\right) ^{-1}\left| c\left( x\right) \right\rangle  \nonumber \\
&&P_{\left\{ 
%TCIMACRO{\func{range}}%
%BeginExpansion
\mathop{\rm range}%
%EndExpansion
{\cal A}\left( x\right) \right\} ^{c}}\left| c\left( x\right) \right\rangle
=0
\end{eqnarray}

\section{Antisymmetric Operators}

If the operator Let ${\cal A}$ be a compact antisymmetric operator $\in 
{\cal B}_{0}\left( V\right) $, where $V$ is a Hilbert space then it has the
canonical expansion 
\begin{equation}
{\cal A}=\sum_{1\leq i\leq s}\alpha _{i}\left\{ \left| s_{i}\right\rangle
\left\langle s_{i+s}^{d}\right| -\left| s_{i+s}\right\rangle \left\langle
s_{i}^{d}\right| \right\}  \label{antiexp}
\end{equation}%
If such an operator is also real, ${\cal A}={\cal A}^{\ast }$, it can be
converted into a compact hermitian one, ${\cal H}$, with the properties 
\begin{eqnarray}
{\cal H} &=&i{\cal A}  \nonumber \\
{\cal H}^{\dagger } &=&-i{\cal A}^{t}=i{\cal A}={\cal H}  \nonumber \\
{\cal H}^{\ast } &=&-i{\cal A}=i{\cal A}^{t}={\cal H}^{t}
\end{eqnarray}%
and 
\begin{equation}
{\cal H}=i{\cal A}=\sum_{1\leq i\leq s}\alpha _{i}\left\{ \left|
v_{i}\right\rangle \left\langle v_{i}\right| -\left| v_{i}^{\ast
}\right\rangle \left\langle v_{i}^{\ast }\right| \right\} ;\quad \alpha
_{i}\in {\Bbb R}
\end{equation}%
which is seen by observing that 
\begin{equation}
{\cal H}v=\alpha v\Rightarrow {\cal H}^{\ast }v^{\ast }=\alpha v^{\ast
}\Rightarrow {\cal H}^{t}v^{\ast }=\alpha v^{\ast }\Rightarrow {\cal H}%
v^{\ast }=-\alpha v^{\ast }
\end{equation}%
Thus 
\begin{equation}
{\cal A}=-\sum_{1\leq j\leq s}i\alpha _{j}\left\{ \left| v_{j}\right\rangle
\left\langle v_{j}\right| -\left| v_{j}^{\ast }\right\rangle \left\langle
v_{j}^{\ast }\right| \right\} ;\quad \alpha _{j}\in {\Bbb R}
\end{equation}%
We can define 
\begin{eqnarray}
\left| s_{i}\right\rangle &=&\left| v_{j}+v_{j}^{\ast }\right\rangle 
\nonumber \\
\left| s_{i+s}\right\rangle &=&i\left| v_{j}-v_{j}^{\ast }\right\rangle
\end{eqnarray}%
then 
\begin{eqnarray}
{\cal A}\left| v_{j}+v_{j}^{\ast }\right\rangle &=&-i\alpha _{j}\left|
v_{j}\right\rangle +i\alpha \left| v_{j}^{\ast }\right\rangle =-\alpha
i\left| v_{j}-v_{j}^{\ast }\right\rangle  \nonumber \\
{\cal A}i\left| v_{j}-v_{j}^{\ast }\right\rangle &=&\alpha _{j}\left|
v_{j}\right\rangle +\alpha _{j}\left| v_{j}^{\ast }\right\rangle =\alpha
_{j}\left| v_{j}+v_{j}^{\ast }\right\rangle
\end{eqnarray}%
thus 
\begin{equation}
{\cal A=}\sum_{1\leq i\leq s}\alpha _{i}\left\{ \left| s_{i}\right\rangle
\left\langle s_{i+s}^{d}\right| -\left| s_{i+s}\right\rangle \left\langle
s_{i}^{d}\right| \right\}
\end{equation}%
which justifies the expansion in eq. $\left( \ref{antiexp}\right) $.
